LLM-based agents increasingly depend on correct tool selection for downstream success, yet
the internal mechanisms that drive this behavior remain poorly understood.
Existing work measures selection accuracy but cannot explain or predict failures.
We present the first empirical mechanistic characterization of tool selection in \LLMs,
combining behavioral analysis, semantic similarity probing, causal perturbation, and
layer-wise representational probing across four complementary experiments.

Using a custom dataset of 55 tool selection scenarios spanning five categories
(weather, calculation, web search, code execution, translation),
we find that cosine similarity between query and tool description embeddings predicts
40\% of selections---64\% better than chance---yet falls short of fully explaining model
behavior.  In controlled rotation experiments, selections remain stable across 90\% of
tool orderings, demonstrating that semantic content, not position, primarily drives
selection when descriptions are informative.  Causal perturbation confirms this:
swapping tool descriptions causes the model to follow the description over the tool name
in 67\% of cases.

Mechanistically, we probe \gpttwo{}'s representations layer by layer and uncover a
striking dissociation.  Tool \emph{category} is perfectly decodable from every layer
(100\% probe accuracy from layer~1), suggesting it is largely encoded at the
token--vocabulary level.  In contrast, information about \emph{which specific tool
position} is semantically correct rises progressively from 35\% at layer~1 to 68\% at
layer~9 (chance~=~25\%), peaking before the final layer.  These findings establish that
semantic matching is a non-trivial computation distributed across the middle-to-late
transformer layers---distinct from coarse category routing---with direct implications for
debugging tool selection failures in deployed agents.
